\label{sec:method}
In this section, we present the Restart sampling algorithm, which incorporates stochasticity during sampling while enabling fast generation. We introduce the algorithm in Sec~\ref{sec:restart-method}, followed by a theoretical analysis in Sec~\ref{sec:ana}. Our analysis shows that Restart achieves a better Wasserstein upper bound compared to those of SDE and ODE in Theorem~\ref{thm1_informal} due to greater contraction effects. %, demonstrating its capability in better generation quality and fast sampling.

\subsection{Method}
\label{sec:restart-method}

In the Restart algorithm, simulation performs a few repeated back-and-forth steps within a pre-defined time interval $[\tmin,\tmax] \subset [0,T]$, as depicted in Figure~\ref{fig:main-demo}. This interval is embedded into the simulation of the original backward ODE referred to as the \textit{main backward process}, which runs from $T$ to $0$. In addition, we refer to the backward process within the Restart interval $[\tmin,\tmax]$ as the \textit{Restart backward process}, to distinguish it from the main backward process.

Starting with samples at time $\tmin$, which are generated by following the main backward process, the Restart algorithm adds a large noise to transit the samples from $\tmin$ to $\tmax$ with the help of the forward process. The forward process does not require any evaluation of the neural network ${s_\theta(x, t)}$, as it is generally defined by an analytical perturbation kernel capable of transporting distributions from $\tmin$ to $\tmax$. For instance, in the case of diffusion models, the perturbation kernel is $\gN(\bm{0}, (\sigma(\tmax)^2 - \sigma(\tmin)^2)\bm{I}_{d\times d})$. 
The added noise in this step induces a more significant contraction compared to the small, interleaved noise in SDE. 
The step acts as if partially restarting the main backward process by increasing the time. Following this step, Restart simulates the backward ODE from $\tmax$ back to $\tmin$ using the neural network predictions as in regular ODE. We repeat these forward-backward steps within $[\tmin,\tmax]$ interval $K$ times in order to further derive the benefit from contraction. Specifically, the forward and backward processes in the $i^{\textrm{th}}$ iteration ($i \in \{0, \dots, K-1\}$) proceed as follows:
\begin{align*}
(\textrm{Restart forward process}) \qquad x^{i+1}_{\tmax} &= x^i_{\tmin} + \varepsilon_{\tmin \to \tmax} \numberthis \label{eq:forward}\\
(\textrm{Restart backward process}) \qquad x^{i+1}_{\tmin} &= \bODE(x^{i+1}_{\tmax}, \tmax \to \tmin)\numberthis \label{eq:backward}
\end{align*}
where the initial $x^0_{\tmin}$ is obtained by simulating the ODE until $\tmin$: $x^{0}_{\tmin} = \bODE(x_{T}, T \to \tmin)$, and the noise $\varepsilon_{\tmin \to \tmax}$ is sampled from the corresponding perturbation kernel from $\tmin$ to $\tmax$. 
% We denote the repetition of the above restart process as the function $\textrm{Restart}_{\tmax \to \tmin, K}$. 
The Restart algorithm not only adds substantial noise in the Restart forward process~(\Eqref{eq:forward}), but also separates the stochasticity from the ODE, leading to a greater contraction effect, which we will demonstrate theoretically in the next subsection. For example, we set $[\tmin, \tmax]=[ 0.05, 0.3]$ for the VP model~\cite{Karras2022ElucidatingTD} on CIFAR-10. Repetitive use of the forward noise effectively mitigates errors accumulated from the preceding simulation up until $\tmax$. Furthermore, the Restart algorithm does not suffer from large discretization errors as it is mainly built from following the ODE in the Restart backward process~(\Eqref{eq:backward}). The effect is that the Restart algorithm is able to reduce the total sampling errors even in the small NFE regime. Detailed pseudocode for the Restart sampling process can be found in Algorithm~\ref{alg:restart}, Appendix~\ref{app:restart-alg}.


\subsection{Analysis}
 \label{sec:ana}

We provide a theoretical analysis of the Restart algorithm under the same setting as Theorem~\ref{thm1_informal}. In particular, we prove the following theorem, which shows that Restart achieves a much smaller contracted error in the Wasserstein upper bound than SDE~(Theorem~\ref{thm1_informal}), thanks to the separation of the noise from the drift, as well as the large added noise in the Restart forward process~(\Eqref{eq:forward}). The repetition of the Restart cycle $K$ times further leads to a enhanced reduction in the initial accumulated error. We denote the intermediate distribution in the $i^{\textrm{th}}$ Restart iteration, following the discretized trajectories and the learned field $s_\theta$, as {\footnotesize $p^{\bRestart(i)}_{t \in [\tmin, \tmax]}$}.

% suggesting faster mixing of the model distribution and the data distribution. 

\begin{theorem}[Informal]
\label{thm2_informal}
Under the same setting of Theorem~\ref{thm1_informal}, 
% let the restarting algorithm be
% \begin{align*}
% \bx^K(\tmin) = \bRestart_{\tmax \to \tmin, K}(\bx(\tmax)).
% \end{align*}
% let $p^{\bRestart, i}_{\tmin}$ denote the distributions for $\bx^i(\tmin)$ and $L_2$ be a uniform bound on $\frac{\lrn{s_{\theta}(x,t) - s_{\theta}(y,t)}}{\lrn{x-y}}$ for all $x,t$. 
assume $K \leq \frac{C}{L_2 \lrp{\tmax - \tmin}}$ for some universal constant $C$. Then
{\footnotesize
\begin{align*}
   \underbrace{W_1(p^{\bRestart(K)}_{\tmin}, p_{\tmin})}_{\textrm{total error}} \leq& \underbrace{B \cdot \lrp{1 - \lambda}^K TV(p^{\bRestart(0)}_{\tmax}, p_{\tmax})}_{\textrm{upper bound on contracted error}} + \underbrace{(K+1) \cdot O\lrp{\delta + \eapprox} \lrp{\tmax - \tmin}}_{\textrm{upper bound on additional sampling error}}
\end{align*}}
where $\lambda<1$ is the same contraction factor as Theorem~\ref{thm1_informal}. $O()$ hides polynomial dependency on various Lipschitz constants, dimension.
\end{theorem}
\vspace{-6pt}
\begin{proofs}
To bound the total error, we introduce an auxiliary process {\footnotesize $q^{\bRestart(i)}_{t \in [\tmin, \tmax]}$}, which initiates from true distribution $p_{\tmax}$ and performs the Restart iterations. This process differs from {\footnotesize ${p}^{\bRestart(i)}_{t \in [\tmin, \tmax]}$} only in its initial distribution at $\tmax$~($p_{\tmax}$ versus {\footnotesize $p^{\bRestart(0)}_{\tmax}$}). We bound the total error by the following triangular inequality:

{\footnotesize
\begin{align*}
    \underbrace{W_1(p^{\bRestart(K)}_{\tmin}, p_{\tmin})}_{\textrm{total error}} \le \underbrace{W_1(p^{\bRestart(K)}_{\tmin}, q^{\bRestart(K)}_{\tmin})}_{\textrm{contracted error}} + \underbrace{W_1(q^{\bRestart(K)}_{\tmin}, p_{\tmin})}_{\textrm{additional sampling error}}
\end{align*}}

To bound the contracted error, we construct a careful coupling process between two individual trajectories sampled from {\footnotesize $p^{\bRestart(i)}_{\tmin}$} and {\footnotesize $q^{\bRestart(i)}_{\tmin}, i=0, \dots, K-1$}. Before these two trajectories converge, the Gaussian noise added in each Restart iteration is chosen to maximize the probability of the two trajectories mapping to an identical point, thereby maximizing the mixing rate in TV. After converging, the two processes evolve under the same Gaussian noise, and will stay converged as their drifts are the same. Lastly, we convert the TV bound to $W_1$ bound by multiplying $B$. The bound on the additional sampling error echoes the ODE analysis in Theorem \ref{thm1_informal}: since the noise-injection and ODE-simulation stages are separate, we do not incur the higher discretization error of SDE.
    \vspace{-6pt}
\end{proofs}
We defer the formal version and proof of Theorem \ref{thm2_informal} to Appendix \ref{a:main_result}. The first term in RHS bounds the contraction on the initial error at time $\tmax$ and the second term reflects the additional sampling error of ODE accumulated across repeated Restart iterations. Comparing the Wasserstein upper bound of SDE and ODE in Theorem~\ref{thm1_informal}, we make the following three observations: \textit{(1)} Each Restart iteration has a smaller contraction factor $1-\lambda$ compared to the one in SDE, since Restart separates the large additive noise~(\Eqref{eq:forward}) from the ODE~(\Eqref{eq:backward}). \textit{(2)} Restart backward process~(\Eqref{eq:backward}) has the same order of discretization error $O(\delta)$ as the ODE, compared to $O(\sqrt{\delta})$ in SDE. Hence, the Restart allows for small NFE due to ODE-level discretization error. \textit{(3)} The contracted error further diminishes exponentially with the number of repetitions $K$ though the additional error increases linearly with $K$. It suggests that there is a sweet spot of $K$ that strikes a balance between reducing the initial error and increasing additional sampling error. Ideally, one should pick a larger $K$ when the initial error at time $\tmax$ greatly outweigh the incurred error in the repetitive backward process from $\tmax$ to $\tmin$. We provide empirical evidences in Sec~\ref{sec:exp-standard}.

While Theorem~\ref{thm1_informal} and Theorem~\ref{thm2_informal} compare the upper bounds on errors of different methods, we provide empirical validation in Section~\ref{sec:exp-verify} by directly calculating these errors, showing that the Restart algorithm indeed yields a smaller total error due to its superior contraction effects. The main goal of Theorem~\ref{thm1_informal} and Theorem~\ref{thm2_informal} is to study how the already accumulated error changes using different samplers, and to understand their ability to self-correct the error by stochasticity. In essence, these theorems differentiate samplers based on their performance post-error accumulation. For example, by tracking the change of accumulated error, Theorem 1 shed light on the distinct "winning regions" of ODE and SDE: ODE samplers have smaller discretization error and hence excel at the small NFE regime. In contrast, SDE performs better in large NFE regime where the discretization error is negligible and its capacity to contract accumulated errors comes to the fore.



\subsection{Practical considerations}
\label{sec:practical}

The Restart algorithm offers several degrees of freedom, including the time interval $[\tmin, \tmax]$ and the number of restart iterations $K$. Here we provide a general recipe of parameter selection for practitioners, taking into account factors such as the complexity of the generative modeling tasks and the capacity of the network. Additionally, we discuss a stratified, multi-level Restart approach that further aids in reducing simulation errors along the whole trajectories for more challenging tasks.

\textbf{Where to Restart? } Theorem~\ref{thm2_informal} shows that the Restart algorithm effectively reduces the accumulated error at time $\tmax$ by a contraction factor in the Wasserstein upper bound. These theoretical findings inspire us to position the Restart interval $[\tmin, \tmax]$ towards the end of the main backward process, where the accumulated error is more substantial. In addition, our empirical observations suggest that a larger time interval $\tmax {-} \tmin$ is more beneficial for weaker/smaller architectures or more challenging datasets. Even though a larger time interval increases the additional sampling error, the benefits of the contraction significantly outweighs the downside, consistent with our theoretical predictions. We leave the development of principled approaches for optimal time interval selection for future works.

\textbf{Multi-level Restart } 
For challenging tasks that yield significant approximation errors, the backward trajectories may diverge substantially from the ground truth even at early stage. To prevent the ODE simulation from quickly deviating from the true trajectory, we propose implementing multiple Restart intervals in the backward process, alongside the interval placed towards the end. Empirically, we observe that a $1$-level Restart is sufficient for CIFAR-10, while for more challenging datasets such as ImageNet~\cite{Deng2009ImageNetAL}, a multi-level Restart results in enhanced performance~\cite{Deng2009ImageNetAL}.